\section{Related Work}
\label{sec:rw}

Infrastructure as Code (IaC) has become a cornerstone of modern DevOps practices, enabling practitioners to define, provision, and manage cloud infrastructure through declarative or imperative scripts. However, the automation and scale that IaC provides also amplify the risk of security misconfigurations, which can persist for extended periods and affect multiple deployed systems simultaneously. This literature review surveys the state of the art in IaC security smell detection, examining empirical studies of security weaknesses, static analysis tools, machine learning approaches, transformer-based code models, and hybrid detection systems that combine rule-based and learning-based techniques.

\subsection{Security Smells in Infrastructure as Code: Empirical Studies}

The concept of \emph{security smells}—recurring coding patterns indicative of security weaknesses—was systematically introduced to the IaC domain by Rahman, Parnin, and Williams \cite{Rahman19}, who conducted a foundational empirical study on Puppet scripts. Analyzing 15,232 scripts from 293 open-source repositories, they identified seven security smell categories: hard-coded secrets, use of HTTP without TLS, suspicious comments, use of weak cryptographic algorithms, invalid IP address binding, empty passwords, and admin-by-default configurations. Their static analysis tool, SLIC (Security Linter for Infrastructure as Code Scripts), detected 21,201 security smell occurrences, including 1,326 hard-coded passwords. Of 1,000 submitted bug reports, 67 were accepted by development teams, validating the relevance of these smells in practice. Notably, some hard-coded secrets persisted for as long as 98 months, with a median lifetime of 20 months, demonstrating the long-term security risks inherent in IaC scripts.

Building on this work, Rahman et al. \cite{Rahman20} extended their study to Ansible and Chef scripts in a replication study. Applying qualitative analysis to 1,956 scripts, they identified two additional security smells not reported in prior work: missing default cases in conditional statements and absence of integrity checks for downloaded artifacts. Their extended tool, SLAC (Security Linter for Ansible and Chef Scripts), analyzed 50,323 scripts from 813 repositories, identifying 46,600 security smell occurrences, including 7,849 hard-coded passwords. With 65 of 94 responded bug reports accepted by practitioners, the study confirmed that security smells are prevalent across multiple IaC languages and remain a persistent concern in DevOps workflows.

Recent empirical work by Silva et al. \cite{Silva24} assessed the adoption of security policies by developers across cloud providers, analyzing 812 Terraform repositories with Checkov and tfsec. Their findings revealed that while access policies were widely adopted, encryption-at-rest policies were consistently neglected across AWS, Azure, and Google Cloud, pointing to systematic gaps in practitioner awareness or tooling support for certain security best practices.

\subsection{Static Analysis Tools for Infrastructure as Code}

Traditional static analysis tools for IaC rely on fixed rule sets to detect syntactically enumerable misconfigurations. Tools such as Checkov \cite{Checkov}, tfsec \cite{Tfsec}, TFLint, and Terrascan have become industry standards, each offering rule-based detection for specific IaC languages. A comparative analysis by Kumar et al. \cite{Kumar24} evaluated these tools on key features and performance metrics, finding that while they excel at detecting permissive file modes, unencrypted protocols, and world-open network rules, they struggle with misconfigurations that depend on contextual or architectural factors rather than fixed syntactic patterns.

To address the limitations of single-language tools, Saavedra and Ferreira \cite{Saavedra23} introduced GLITCH, a technology-agnostic framework for polyglot code smell detection. GLITCH transforms IaC scripts written in Ansible, Chef, Docker, Puppet, or Terraform into an intermediate representation, enabling unified detection of nine security smells and nine design-and-implementation smells. Empirical studies with GLITCH demonstrated higher precision and recall than state-of-the-art tools, reducing the effort required to write smell analyses for multiple IaC technologies. Gonçalves \cite{Goncalves23} extended GLITCH to support Terraform, defining a taxonomy of 28 security smells and comparing its performance against tfsec, showing improved detection capabilities for context-dependent misconfigurations.

Advancing beyond pattern-based detection, Opdebeeck and Zerouali \cite{Opdebeeck23} introduced GASEL, a novel Ansible security smell detector that uses graph queries on program dependence graphs (PDGs) to detect seven security smells. Unlike existing approaches that consider individual files (akin to intraprocedural analysis), GASEL performs whole-project, inter-procedural analysis and leverages control-flow and data-flow information to guide smell detection. Evaluated on an oracle of 243 real-world security smells, GASEL substantially outperformed state-of-the-art detectors in both precision and recall. Their analysis of over 15,000 Ansible scripts revealed that over 55\% of security smells contain data-flow indirection and over 32\% require control-flow awareness, demonstrating the value of program analysis techniques beyond simple syntactic matching.

Complementing detection-focused approaches, TerrARA \cite{Chowdhury24} introduced automated security threat modeling for Terraform configurations. TerrARA automatically constructs enriched Data Flow Diagrams (DFDs) from Terraform files for AWS resources, encodes cloud computing threat patterns, and utilizes the SPARTA threat modeling engine to identify architectural security flaws such as broken access control or insecure storage. This approach addresses a gap in existing tools, which predominantly focus on implementation bugs (e.g., hardcoded secrets) rather than design-level security weaknesses.

\subsection{Machine Learning for Vulnerability and Misconfiguration Detection}

Machine learning has emerged as a complementary approach to rule-based static analysis, particularly for identifying complex or subtle security vulnerabilities that evade fixed patterns. Early work in ML-based vulnerability detection relied on traditional vectorization techniques. Li et al. \cite{Li20} conducted an empirical study comparing deep learning models (BiLSTM, CNN, RNN) with traditional machine learning (SVM, Random Forest, Logistic Regression) for vulnerability detection, finding that deep learning models achieved better overall performance, with BiLSTM being the most effective. However, for traditional ML approaches, CountVectorizer outperformed other vectorization methods, and features extracted from vulnerability datasets closely resembled human expert knowledge.

Medeiros et al. \cite{Medeiros20} demonstrated that TF-IDF combined with decision trees, SVM, or logistic regression could achieve high accuracy (0.97) on the OWASP Benchmark dataset for Java vulnerability prediction, establishing TF-IDF-based approaches as a strong baseline for security detection tasks. Supporting this finding, empirical work by Kovalenko et al. \cite{Kovalenko23} showed that "old-fashioned" TF-IDF with Binary Relevance remains competitive with deep pre-trained models for vulnerability type identification, requiring 40× less training time and 5× less inference time while achieving comparable F1 scores.

Recent work by Nguyen et al. \cite{Nguyen24} investigated the contributing factors for ML-based vulnerability detection across 17 real-world projects, finding that bag-of-words embedding combined with Random Forest consistently increased detection accuracy by approximately 4\% compared to other combinations, though they also observed limitations in transferring vulnerability signatures across domains.

\subsection{Transformer-Based Code Models}

The advent of transformer-based language models has revolutionized natural language processing and, more recently, code analysis. The foundational BERT architecture \cite{Devlin19} introduced bidirectional self-attention, where each token attends to all other tokens in the sequence, enabling rich contextual representations. Building on BERT, Feng et al. \cite{Feng20} introduced CodeBERT, a bimodal pre-trained model for programming and natural language. CodeBERT uses the RoBERTa-base architecture (125M parameters) and is trained with masked language modeling and replaced token detection on NL-PL pairs, enabling it to capture semantic connections between code and natural language comments. CodeBERT has achieved state-of-the-art performance on natural language code search and code documentation generation tasks.

For processing long IaC scripts, Beltagy, Peters, and Cohan \cite{Beltagy20} introduced Longformer, a transformer variant with an attention mechanism that scales linearly rather than quadratically with sequence length. Longformer combines local windowed attention (analogous to CNNs) with task-specific global attention, making it suitable for documents of thousands of tokens. This architecture addresses the fundamental limitation of standard transformers, which cannot efficiently process long Puppet or Terraform scripts without truncation or chunking.

Several recent studies have applied transformer-based models to vulnerability detection. Huang et al. \cite{Huang22} introduced BBVD, a BERT-based method for vulnerability detection. Zhang et al. \cite{Zhang24} explored adversarial attacks against CodeBERT, revealing vulnerabilities in code models themselves. More advanced architectures have emerged: Li et al. \cite{Li23} proposed XGV-BERT, combining CodeBERT with Graph Convolutional Networks (GCNs) to jointly leverage code embeddings and program dependency graphs for C/C++ vulnerability detection. Chen et al. \cite{Chen24} introduced VulD-CodeBERT, which fine-tunes CodeBERT and adds a BiLSTM network with attention mechanisms to mine local key features of vulnerabilities, combined with global features from CodeBERT's pooler output. Mim et al. \cite{Mim24} proposed VulBertCNN, which applies CodeBERT embeddings to nodes in Program Dependency Graphs (PDGs) and uses centrality measures with CNNs, achieving 21.7\% higher accuracy and 26.3\% higher F1 scores than prior state-of-the-art on SARD and Big-Vul datasets.

Tang et al. \cite{Tang24} conducted a comparative empirical study with Defect-Scanner, finding that CodeBERT-based models achieved approximately 90\% accuracy for vulnerability detection in C/C++ programs, substantially outperforming Word2Vec-based approaches. Kumar et al. \cite{Kumar24b} evaluated 540 different models combining CodeBERT embeddings with five feature selection techniques and fifteen classification algorithms, finding that ensembles of decision trees (Random Forest and Extra Trees) performed best, and that CodeBERT embeddings effectively captured vulnerability characteristics even for lightweight classifiers.

\subsection{Baseline: Semantics-Aware Detection with Code and Language Models}

The baseline for this project is War et al. \cite{War25}, ``Detection of Security Smells in IaC Scripts through Semantics-Aware Code and Language Processing'' (University of Luxembourg, arXiv:2509.18790, September 2025). \emph{This is cited as a preprint; its peer-review status could not be confirmed at the time of writing, unlike the baselines used elsewhere in this thesis series.} The paper fine-tunes CodeBERT for Ansible scripts (which are rich in natural-language comments and task descriptions) and LongFormer for long Puppet scripts to classify IaC snippets as misconfigured or safe. Their approach substantially outperforms prior TF-IDF/Bag-of-Words + Random Forest baselines (e.g., Ansible F1 0.90 vs. 0.77) and performs competitively with much larger general-purpose LLMs such as GPT-4, Claude 3.5, and LLaMA.

\subsection{The Gap This Project Targets}

The paper's own ablation study (their RQ2) is the source of the gap this project targets. Removing natural-language context—stripping comments from Ansible snippets or truncating surrounding code for Puppet snippets—causes precision to collapse sharply while recall holds steady or rises: CodeBERT on Ansible drops from precision 0.92 to 0.46 (F1 0.90 to 0.58); LongFormer on Puppet drops from 0.87 to 0.55 (F1 0.79 to 0.70). Their own words: ``without sufficient semantic context... models over-flag and lose precision, undermining practical usability.'' This is a real deployment concern: much real-world IaC is sparsely commented, unlike the curated, comment-rich datasets used to train and validate the paper's models. No fix for this precision collapse is proposed or evaluated in the baseline paper.

\subsection{Hybrid Detection Approaches: Combining Rules and Machine Learning}

Combining rule-based static analysis with machine learning classifiers has emerged as a promising approach for addressing the complementary weaknesses of each technique. Rule-based systems offer perfect precision for enumerable patterns but lack coverage for complex or context-dependent vulnerabilities, while ML models offer broader coverage but can suffer from high false-positive rates.

Schuster et al. \cite{Schuster23} introduced Fluffy, a bimodal taint analysis that combines CodeQL's static data-flow analysis with machine learning to filter unexpected security flows. The key insight is that static analysis reasons about \emph{where} data flows, while machine learning probabilistically determines \emph{whether} the flow is problematic based on natural language information such as API names. Evaluated on five common vulnerability types across 250,000 JavaScript projects, Fluffy achieved an F1 score of 0.85 or higher on four of them, demonstrating the effectiveness of letting ML models predict from natural language information whether a taint flow is expected or unexpected.

Strittmatter et al. \cite{Strittmatter24} proposed Dev-Assist, an IntelliJ IDEA plugin that uses multi-label machine learning to detect security-relevant methods, considering dependencies among labels. Unlike binary relevance approaches, their multi-label approach generalizes better and achieves higher F1 scores. The plugin can automatically generate configurations for Static Application Security Testing (SAST) tools, run static analysis, and display results within the IDE, reducing the manual effort required to configure and use security analysis tools.

For IaC specifically, recent work by Bansal et al. \cite{Bansal26} applied semantic transformers to Terraform misconfiguration detection, leveraging DistilBERT for sequence classification. Their approach achieved an F1 score of 0.7598 and recall of 0.8700 on a balanced dataset of 2,000 samples, outperforming traditional TF-IDF with Logistic Regression in recall and overall detection capability, demonstrating the ability of transformer models to capture security patterns beyond syntactic rules.

\subsection{Positioning of This Work}

This project builds directly on the precision-collapse finding from War et al. \cite{War25} and addresses it with a hybrid detector that combines comment-independent rule-based detection (for enumerable patterns) with a stricter-threshold ML classifier (for patterns that require learned representations). Unlike Fluffy \cite{Schuster23}, which uses ML to \emph{filter} static analysis results, this hybrid approach uses rules to \emph{augment} ML predictions, trading recall for precision recovery in comment-stripped scenarios. The approach is evaluated on a synthetic IaC benchmark explicitly partitioned into "easy" (rule-detectable) and "hard" (comment-dependent) patterns, making the limitations of the hybrid approach explicit: it can only recover precision for the subset of misconfigurations that have enumerable, comment-independent signatures. This honest framing—demonstrating both what the hybrid \emph{can} fix (precision on easy patterns) and what it \emph{cannot} fix (recall on hard patterns that genuinely require context)—is a key contribution to understanding the boundaries of hybrid detection systems for IaC security.
